\documentclass[11pt]{book}

\usepackage[utf8]{inputenc}
\usepackage[T1]{fontenc}
\usepackage{amsmath,amssymb,amsthm,mathtools}
\usepackage{geometry}
\usepackage{hyperref}
\usepackage{booktabs}
\usepackage{longtable}
\usepackage{array}
\geometry{margin=1in}

\newtheorem{theorem}{Theorem}[chapter]
\newtheorem{lemma}[theorem]{Lemma}
\newtheorem{proposition}[theorem]{Proposition}
\newtheorem{corollary}[theorem]{Corollary}
\theoremstyle{definition}
\newtheorem{definition}[theorem]{Definition}
\newtheorem{remark}[theorem]{Remark}
\newcolumntype{L}[1]{>{\raggedright\arraybackslash}p{#1}}

\title{Contradiction Fields: CRT Rigidity, Prime Matrix Geometry, and Zeta-Zero Escape}
\author{Project manuscript}
\date{\today}

\begin{document}
\maketitle

\chapter*{Abstract}
This monograph draft unifies two contradiction-field research programs.  The first studies prime existence in CRT prime matrices through row, column, diagonal, anchor, and rough-composite rigidity.  The second studies an off-critical zero of the Riemann zeta function through a global contradiction-field ledger for prime anomalies.  The common method is to convert a hypothetical counterexample into a structured covering field, decompose it into local rigidity layers, and exclude all free escape channels by capacity, periodicity, and no-cycle bookkeeping.

This is a synthesis and dependency manuscript.  It does not assert that the Riemann Hypothesis or the prime-matrix row/column theorem has already been accepted as a final unconditional theorem.  Each claim is marked by status in the accompanying status table: proved in manuscript, reduced to a structured input, external, computationally certified, or requiring independent referee verification.

\tableofcontents

\chapter{Guiding Principle: The Contradiction Field}

A contradiction field is a structured ledger attached to a hypothetical counterexample.  Instead of estimating a single exceptional set, it records every possible explanation for the counterexample as a named channel.  A proof closes when every channel is either absorbed, forced to descend in a finite potential, transferred with source deletion, or contradicts a global capacity lower/upper bound.

The common rigidities used throughout this manuscript are:
\begin{enumerate}
  \item CRT skeletons and nonzero-residue balance;
  \item local short-window non-reuse of large prime factors;
  \item small-prime shell migration and diagonal locking;
  \item rough-composite and double-anchor decompositions;
  \item capacity ledgers after baseline subtraction;
  \item no-cycle potentials for internal escape mechanisms.
\end{enumerate}

\chapter{Prime-Matrix Contradiction Fields}

Let $P$ be an odd prime and write the $P\times P$ prime matrix entries as
\[
  n(r,c)=(r-1)P+c,\qquad 1\le r,c\le P.
\]
The row and column programs attempt to show that every nontrivial row or column contains a prime.  The current formal appendix reduces a row or non-$P$ column counterexample to a structured bad configuration after removing small-factor locks and tail anchors.

\begin{theorem}[A/B reduction, recorded form]
A nontrivial full-composite row or column, after small-factor and tail-anchor peeling, induces a Structured-EHPD bad configuration satisfying CRT balance, large-factor non-reuse, and anchor-capacity constraints.
\end{theorem}
\begin{proof}[Reference]
See \texttt{docs/row-column-reduction-formal-appendix.md}.  This theorem is a reduction statement; it does not by itself prove the Structured-EHPD exclusion theorem.
\end{proof}


\section{Two-Point Rough-Pair Propositions}

A natural next strengthening asks for two simultaneous nonzero-residue constraints in the back half of the prime matrix.  Fix an even integer $w>0$ and ask whether there is an entry $x$ in the back half of the $P\times P$ matrix such that both $x$ and $x-w$ are nonzero modulo every prime $p\le P$, or equivalently both $x$ and $w-x$ are nonzero modulo every prime $p\le P$.

\begin{proposition}[Equivalence and strength of the two-point rough-pair problem]
For fixed even $w>0$, the two conditions
\[
  p\nmid x(x-w)\quad(p\le P)
\]
and
\[
  p\nmid x(w-x)\quad(p\le P)
\]
are identical.  In the region $1<x\le P^2$, any integer nonzero modulo every prime $p\le P$ is prime.  Consequently, a proof of this two-point assertion for every sufficiently large $P$ would imply infinitely many prime pairs of gap $w$; in particular $w=2$ would imply the twin-prime conjecture.
\end{proposition}

\begin{proof}
For each $p\le P$, the condition $p\nmid x-w$ is $x\not\equiv w\pmod p$, and the condition $p\nmid w-x$ is the same congruence.  Both formulations therefore exclude exactly the two residue classes $0$ and $w$ modulo each $p$, with possible collision when $p\mid w$.

If $1<x\le P^2$ is nonzero modulo every prime $p\le P$ and is composite, then it has a prime factor $q\le \sqrt{x}\le P$, contradicting the nonzero-residue condition.  Thus such an $x$ is prime.  If the two-point assertion holds for all sufficiently large $P>w$, then $x$ and $x-w$ are both prime in infinitely many growing back-half matrix intervals, giving infinitely many prime pairs of gap $w$.
\end{proof}

\begin{remark}[Status]
This two-point problem is therefore not a routine corollary of the one-point row/column contradiction field.  It is a prime-pair-level strengthening, naturally expressible as a two-point rough Jacobsthal problem.  In this monograph it is recorded as a future contradiction-field direction, not as an unconditional theorem.
\end{remark}



\section{Two-Point Secondary-Sieve Contradiction Field}

The two-point rough-pair problem admits a more detailed contradiction-field research program.  The current working manuscript is recorded in
\[
  \texttt{docs/monograph/two-point-secondary-sieve-research.md}.
\]
It studies the secondary sieve condition
\[
  x\not\equiv 0,w\pmod p\qquad(p\le P)
\]
inside back-half matrix windows, and tries to exclude a two-point rough-pair counterexample by a hierarchy of local rigidity fields.

\begin{definition}[Two-point secondary-sieve field]
Fix an even integer $w>0$.  A two-point secondary-sieve field is the structured ledger obtained from a hypothetical interval in which every candidate $x$ fails at least one of the two congruence conditions $p\nmid x$ or $p\nmid x-w$ for some prime $p\le P$.  Its named channels are small-prime double locks, top-scale large-factor coverage, short-block zero-mass loss, multiplicative near-crossings, Fourier CRT defects, and smoothing boundary terms.
\end{definition}

\begin{lemma}[Fixed-order local crossing complexity]
In the secondary-sieve manuscript, every fixed-order local crossing test is expanded only through truncated hard-skeleton weights
\[
  d\le R,\qquad R\le(\log P)^{C_R},
\]
with coefficient mass $\sum_{d\le R}|\lambda_d|\le(\log P)^{C_\lambda}$ and short-block shifts $|h|<L\asymp(\log\log P)^2$.  For every fixed order $k$, the effective historical complexity satisfies
\[
  Q_{\mathrm{eff},k}
  \ll_k
  \left(\sum_{d\le R}|\lambda_d|\right)^kR^kL^{O(k)}
  \le(\log P)^{C_k}.
\]
In particular the required $k=2,3$ local crossing tests have polylogarithmic effective complexity.
\end{lemma}

\begin{proof}
After expanding the $k$ truncated weights, only moduli $d_1,\ldots,d_k\le R$ occur.  For fixed new-line primes $p_i$, phases, and block shifts, all conditions are one-variable linear congruences.  By the Chinese remainder theorem the local system is either inconsistent or occupies one residue class modulo a divisor of
\[
  \operatorname{lcm}(d_1,\ldots,d_k)\operatorname{lcm}(p_1,\ldots,p_k).
\]
The primes $p_i$ are the explicit variables being summed in the crossing estimate; they contribute the usual $Y^k$ scale, not hidden historical complexity.  The historical part is bounded by $\operatorname{lcm}(d_1,\ldots,d_k)\le R^k$, the block shifts contribute $L^{O(k)}$ patterns, and the coefficient expansion contributes $\left(\sum|\lambda_d|\right)^k$.  This proves the displayed bound.  The true residual set $U_Y$ is not expanded into a full primorial period.
\end{proof}

\begin{lemma}[Zero-mass and smoothing discharge]
Assume the SC2/diagonal package supplies the moment bounds
\[
  M_1:=\frac{\sum_BX_BR_B}{\sum_BX_B}\le0.45,\qquad
  M_2:=\frac{\sum_BX_B\binom{R_B}{2}}{\sum_BX_B}\ge0.05,
\]
and
\[
  M_3:=\frac{\sum_BX_B\binom{R_B}{3}_+}{\sum_BX_B}\le0.03.
\]
Then the weighted zero-hit mass satisfies
\[
  \frac{\sum_BX_B1_{R_B=0}}{\sum_BX_B}\ge0.57.
\]
Moreover, if the directional-balance endpoint estimate gives
\[
  N_{\mathrm{endpoint}}/N_{\mathrm{total}}\le C_{\mathrm{end}}\Delta+o(1),
\]
then choosing $\Delta\le0.01/C_{\mathrm{end}}$ gives a smooth-to-sharp loss $<0.03$ for all sufficiently large $P$.
\end{lemma}

\begin{proof}
For every integer $R\ge0$,
\[
  1_{R=0}\ge 1-R+\binom R2-\binom R3_+ .
\]
Multiplying by $X_B\ge0$ and summing over blocks gives
\[
  \frac{\sum_BX_B1_{R_B=0}}{\sum_BX_B}
  \ge 1-M_1+M_2-M_3
  \ge 1-0.45+0.05-0.03=0.57.
\]
For smoothing, choose even smooth cutoffs $W_-\le1_{|t|<1}\le W_+$ whose difference is supported in $1-\Delta\le |t|\le1+\Delta$.  The sharp and smoothed counts differ only on this endpoint layer.  The displayed directional-balance bound therefore gives loss at most $C_{\mathrm{end}}\Delta+o(1)$, which is $<0.03$ after the stated choice of $\Delta$ and then taking $P$ large.
\end{proof}

\begin{lemma}[Adaptive true-hit layering dichotomy]
Let $U_Y$ be the current true residual set and define, for each new prime $p>Y$,
\[
  a_p=\frac1{|U_Y|}\#\{x\in U_Y:x\equiv0\text{ or }w\pmod p\}.
\]
Fix $\eta=0.05$.  Either some $p$ has $a_p>\eta$, in which case the configuration is a Single-Prime CRTDefect, or the remaining primes can be greedily partitioned into consecutive layers $\mathcal P_j$ satisfying
\[
  \nu_j^{\mathrm{real}}:=\sum_{p\in\mathcal P_j}a_p\le0.4,
\]
and, except for the final layer, $\nu_j^{\mathrm{real}}>0.4-\eta\ge0.35$.
\end{lemma}

\begin{proof}
If $a_p>\eta$ for some $p$, then a fixed positive proportion of $U_Y$ lies in the two residue classes $0,w$ modulo $p$.  This is exactly the Single-Prime CRTDefect exit, to be excluded by the CRTDefect/Directional Balance part of the SC2 package.  Otherwise all $a_p\le\eta$.  Starting from the smallest remaining prime, add primes to the current layer until adding the next one would make the partial sum exceed $0.4$.  The completed layer has sum at most $0.4$, and because the next summand is at most $\eta$, every non-final completed layer has sum greater than $0.4-\eta$.  This gives the claimed adaptive true-hit decomposition.
\end{proof}

\begin{proposition}[Recorded conditional closure chain for the two-point secondary sieve]
Using the three lemmas above, assume the following remaining input holds with the constants recorded in the secondary-sieve research note:
\begin{enumerate}
  \item the I3-Core SC2/CRTDefect package controls the true-residual two-line near-crossing error, the odd-prime $45^\circ$ main-synchronization peak, Single-Prime CRTDefect exits, the normalized $q=2$ parity skeleton, the moment bounds $M_1,M_2,M_3$, and the endpoint Directional Balance used in the zero-mass and smoothing discharge lemma.
\end{enumerate}
Then the two-point secondary-sieve contradiction field has no remaining top-scale escape channel under the parameter choice
\[
  \alpha_0=0.75,\qquad c_0=1/2,\qquad \nu_j^{\mathrm{real}}\le0.4
\]
for adaptive true-hit layers, subject to the same fixed-$w$ convention used in the research manuscript.
\end{proposition}

\begin{proof}[Proof sketch and status]
The argument is a conditional synthesis, not a final unconditional theorem.  A top-scale covering excess is first decomposed into adaptive true-hit layers; any single-prime concentration is routed to the Single-Prime CRTDefect exit.  The remaining excess is converted by a TCA linear bridge into a positive short-block load.  The block load is forced into the SMC/SC2 near-crossing ledger.  The $45^\circ$ main synchronization is split into large-factor reuse, shared-complement factors, and good synchronized pairs; reuse and shared factors are excluded by short-window rigidity, while good pairs are sent to an auxiliary-prime Fourier CRT-defect.  Directional balance excludes the Fourier defect unless the top-scale energy is already below the capacity threshold.  Finally, the weighted Bonferroni zero-mass estimate supplies enough zero-hit block mass to make the capacity comparison incompatible.

The current audit is archived under \path{docs/archive/monograph-reviews/} and classifies this as a research proposition conditional on the single I3-Core package above.  The fixed-order complexity input, zero-mass/smoothing discharge, and adaptive layering dichotomy have been inlined as the preceding lemmas.  The fact that an unconditional proof would imply fixed-gap prime-pair consequences, including the twin-prime case when $w=2$, is not a logical obstruction.  It only marks the strength of the claim.  The reason this chapter remains conditional is that I3-Core has not yet been inlined as an independently verified unconditional proof.
\end{proof}

\begin{proposition}[Window-compressed conditional variant]
Keep the same fixed even integer $w>0$ and the same remaining I3-Core input as in the preceding proposition.  Let $I$ be any consecutive back-half window of $H(P)$ rows in the $P\times P$ matrix, so that $|I|=H(P)P+O(P)$.  By the fixed-order complexity lemma, write
\[
  Q_{\mathrm{eff}}\le C_Q(\log P)^C.
\]
For a prescribed middle-scale error budget $\varepsilon_{\mathrm{SC2}}>0$, the conditional row threshold supplied by the current ledger is
\[
  H_{\min}^{\mathrm{cond}}(P;\varepsilon_{\mathrm{SC2}})
  =
  \left\lceil
  \frac{C_Q}{\varepsilon_{\mathrm{SC2}}}
  P^{1/2}(\log P)^C
  \right\rceil.
\]
Equivalently, one may use the asymptotic form
\[
  H(P)\ge H_{\min}^{\mathrm{asym}}(P)
  :=
  \left\lceil P^{1/2}(\log P)^{C_*}\right\rceil,
  \qquad C_*>C.
\]
Under either threshold, the same conditional contradiction-field chain gives an element
\[
  x\in I,\qquad p\nmid x(x-w)\quad(p\le P).
\]
When $x-w>1$, both $x$ and $x-w$ are primes.
\end{proposition}

\begin{proof}[Recorded scale check]
The only scale-sensitive term in the recorded chain is the middle/high-scale $K=2$ near-crossing error.  With top layer $Y=P^{3/4}$ and $Q_{\mathrm{eff}}\le(\log P)^C$, replacing the original back-half length $\asymp P^2$ by $|I|\asymp H(P)P$ changes the normalized error to
\[
  \frac{Q_{\mathrm{eff}}Y^2}{|I|}
  \ll
  \frac{P^{1/2}(\log P)^C}{H(P)}.
\]
The displayed lower bound on $H(P)$ makes this at most $\varepsilon_{\mathrm{SC2}}$ in the explicit form, or $o(1)$ in the asymptotic form.  Hence the constants in the TCA, SC2, zero-mass, finite-package, and smoothing ledgers are unchanged.  The conclusion is therefore a conditional window-compressed version of the previous proposition.  It becomes an unconditional prime-pair theorem only if I3-Core is independently proved and accepted.
\end{proof}

\begin{remark}[Referee boundary]
The two-point secondary-sieve chapter is included to document the contradiction-field architecture and its parameter ledger.  The window-compression audit is recorded in \path{docs/monograph/two-point-window-compression-and-unconditionality.md}.  It must not be cited as an unconditional proof of prime pairs, the twin-prime conjecture, or Goldbach-type assertions unless I3-Core is independently proved and accepted.
\end{remark}

\begin{remark}[Current obstruction]
The current hard attack reduces I3-Core to a true-residual correlation problem.  The local CRT and matrix rigidities do not by themselves exclude a model residual set concentrated on a new prime's bad residue class while still satisfying all previously imposed two-residue exclusions and short-window non-reuse constraints.  Thus the remaining task is not another bookkeeping step but a genuine lower-bound/correlation theorem for the true residual set in new prime moduli.
\end{remark}

\chapter{Zeta-Zero Contradiction Fields}

Assume an off-critical zero $\rho=\beta+i\gamma$ with $\beta>1/2$.  The RH contradiction-field manuscript transfers the resulting smooth prime anomaly into a CRT rough-composite ledger, decomposes it into sparse/dense/tail/internal/global exits, and proves that no named escape channel remains under the recorded local theorems and restricted external inputs.

\begin{theorem}[RH contradiction-field synthesis, recorded form]
Under the numbered local theorems, controlled-exit lemmas, threshold hierarchy, and restricted external inputs recorded in \path{paper/rh-proof/rh-contradiction-field.tex}, an off-critical zero has no free escape channel in the contradiction-field event graph.
\end{theorem}
\begin{proof}[Reference]
See the RH manuscript in \path{paper/rh-proof/} and the final audit in \path{docs/}.  The submission warning in the RH manuscript remains in force until independent line-by-line referee verification and final monograph page-number checks are completed.
\end{proof}

\chapter{Unified Dependency Map}

The two programs share the same logic:
\[
  \text{counterexample}\Rightarrow \text{structured covering field}\Rightarrow
  \text{named exits}\Rightarrow \text{capacity/no-cycle contradiction}.
\]
The prime-matrix program is more geometric and finite-CRT in nature; the zeta-zero program is analytic and scale-asymptotic.  The monograph treats both as projections of the same contradiction-field method.

\begin{longtable}{@{}L{0.22\textwidth}L{0.34\textwidth}L{0.34\textwidth}@{}}
\toprule
Layer & Prime-matrix projection & RH projection\\
\midrule
Baseline & CRT nonzero class balance & smoothed CRT candidate baseline\\
Local rigidity & large-factor non-reuse, diagonal locks & sparse/dense/tail routing\\
Capacity & anchor coverage capacity & projection and square-function capacity\\
No-cycle & recursive peeling and shell migration & GEE transfer/no-cycle ledger\\
External inputs & finite verification and Structured-EHPD & explicit formula, KL, Vaaler, BG/Baker, sieve, Vaughan\\
Final status & reduction package, not final unconditional theorem & verification package with submission warning\\
\bottomrule
\end{longtable}


\chapter{Referee Closure Audit}

This chapter records the strict closure standard used for the whole monograph.  A contradiction-field chain is considered closed only if its entrance, decomposition, exits, and global contradiction are all verified under the same load convention.

\begin{definition}[Closure certificate]
A closure certificate for a contradiction field consists of four data:
\begin{enumerate}
  \item an entrance theorem sending every counterexample into the field;
  \item a disjoint or source-deleting decomposition into named exits;
  \item for each exit, an absorption, descent, transfer, external theorem, or finite certificate;
  \item a global comparison showing that final load is simultaneously bounded below and above in incompatible ways.
\end{enumerate}
If any item is replaced by an unverified structural input, the conclusion is conditional on that input.
\end{definition}

\begin{theorem}[Current monograph closure]
The present manuscript closes the audit and dependency structure of the two contradiction-field programs: all known entrances, decompositions, exits, external inputs, finite certificates, and referee obligations are explicitly named.  It does not close the final unconditional theorem status of either program.
\end{theorem}
\begin{proof}
For the prime-matrix program, the A/B appendix gives the entrance reduction from a row or column counterexample to a Structured-EHPD bad configuration, while the D-structure exclusion remains the decisive input requiring independent line-by-line review.  Therefore the chain is closed as a reduction package but not as a final theorem.

For the RH program, the RH manuscript has completed U1--U5: statement wording, AEX acceptance, EXT localization, manuscript wording, BibTeX/PDF compilation, and warning-free LaTeX logs.  The final audit still requires independent referee verification of every local theorem and controlled exit.  Therefore the chain is closed as a verification package but not as a final unconditional RH theorem.

The status table and the referee audit in \path{docs/monograph/} record exactly which pieces are proved, reduced, external, computational, or awaiting review.  Thus the monograph's logical boundary is explicit and closed, while promotion to final theorem status remains conditional on external verification.
\end{proof}


\begin{theorem}[Internal line-by-line audit]
The author-side line-by-line audit has found no new item classified as a mathematical block.  The remaining blocks are referee blocks: Prime Matrix final promotion requires independent acceptance of the D-structure/Tail-log4/finite-verification interface, and RH final promotion requires independent acceptance of all controlled exits in the RH contradiction-field manuscript.
\end{theorem}
\begin{proof}
The internal audit matrix assigns one of four labels to each link.  No link is a math-block.

The terminal entries PM-16 and RH-16 remain referee-blocks.  An author-side check cannot replace independent referee verification.  Therefore the internal audit is complete, but final unconditional theorem promotion is not.
\end{proof}

\begin{remark}[No overclaim principle]
The monograph may state the shared contradiction-field method and the completed audit closures.  It must not state unconditional proofs of RH or of the prime-matrix row/column theorem until the remaining structured inputs have independent acceptance.
\end{remark}

\chapter{Current Status and Honest Boundary}

The purpose of this monograph draft is to make the shared structure visible and auditable.  It is not yet a final unconditional theorem manuscript.  The precise statuses are maintained in \texttt{docs/monograph/claim-status-table.md}.  Any future promotion to final theorem status requires independent verification of every reduced structured input and every controlled exit.

\end{document}
